\chapter{Détection automatique de langue}

\section{Identification}

\begin{center}
\begin{tabular}{ll}
\toprule
Modèle         & \texttt{facebook/fasttext-language-identification} \\
Nb langues     & 217 \\
Statut scoring & \textbf{Implémenté} \\
Fichier        & \texttt{DataPipelines/LanguageDetection.py} \\
\bottomrule
\end{tabular}
\end{center}

\section{Fonctionnalité}

Premier composant de la chaîne : détermine si une traduction est nécessaire.
Déclenché sur les 700 premiers + 1500--2500 caractères du texte. Retourne le
code langue BCP-47 FLORES-200 (ex.\ \texttt{fra\_Latn}, \texttt{ara\_Arab}).

\section{Fonction de scoring de confiance}

\noindent\textit{Implémentation dans
\texttt{LanguageDetection.compute\_confidence()}.}

\medskip
FastText applique un softmax sur 217 classes. Le score de confiance est
directement la probabilité de la langue prédite :

\begin{equation}
  s_{\text{LID}}(x)
  = P(\hat{l} \mid x)
  = \max_{l \in \mathcal{L}} P(l \mid x)
  \;\in (0, 1]
  \label{eq:lid}
\end{equation}

Les top-2 prédictions sont stockées dans les métadonnées pour auditabilité.

\paragraph{Seuils.}

\begin{center}
\begin{tabular}{lll}
\toprule
$s \geq 0.90$ & Langue certaine & Décision automatique fiable \\
$0.60 \leq s < 0.90$ & Langue probable & Logging, surveiller \\
$s < 0.60$ & Détection ambiguë & Ne pas traduire automatiquement \\
\bottomrule
\end{tabular}
\end{center}

\textbf{Métadonnées loguées :} \texttt{detected\_lang},
\texttt{top\_k} (liste des 2 meilleures prédictions avec scores).

\section{Justification scientifique}

FastText, probabilités softmax~\cite{joulin2017bag} ;
cas difficiles (code-switching, dialectes)~\cite{caswell2020language}.

\section{Validation empirique}

Aucun dataset d'évaluation annoté n'est disponible à ce jour. Le
\texttt{compute\_confidence()} est actuellement validé par inspection manuelle
sur des cas réels et des scénarios limites (\textit{smoke testing}).

\noindent\textbf{Métriques à mesurer lors de la constitution du dataset :}
\begin{itemize}
  \item Corrélation entre le score de confiance et la bonne classification ;
  \item Calibration (ECE) : un score de confiance de 0.8 correspond-il à
    environ 80\,\% de réussite sur le dataset ?
  \item Pouvoir discriminant (AUC-ROC) : le score sépare-t-il les prédictions
    correctes des incorrectes ?
\end{itemize}

\noindent\textbf{Statut :} \textit{À implémenter.}

\paragraph{Stress tests à conduire.}
Textes bilingues mélangés, textes très courts ($< 20$ caractères), arabe
dialectal levantin, translittération arabe en latin.
